% Plantilla mínima de los boletines IES (pandoc --template).
% Solo carga lo estrictamente necesario por el contenido; la tipografía y la
% geometría las decide LaTeX (tamaño carta, 12pt, márgenes por defecto).

\documentclass[12pt,letterpaper]{article}

% Codificación y acentos (pdflatex)
\usepackage[T1]{fontenc}
\usepackage[utf8]{inputenc}

% Fuente escalable (Latin Modern); evita las bitmap EC de METAFONT
\usepackage{lmodern}

% Español: títulos y pies ("Figura"), guionado
\usepackage[spanish]{babel}

% Secciones sin numeración automática (equivalente a \section*{...}): solo se
% muestran los números que el markdown escribe a mano ("1.", "2.", "A.", ...).
\setcounter{secnumdepth}{-1}

% Los títulos nunca quedan huérfanos al pie de página: reserva espacio antes
% de cada sección y, si no cabe el título + unas líneas, salta a la página
% siguiente. Necesario sobre todo cuando el título precede a una tabla longtable.
\usepackage{needspace}
\usepackage{etoolbox}
\pretocmd{\section}{\needspace{3\baselineskip}}{}{}
\pretocmd{\subsection}{\needspace{3\baselineskip}}{}{}
\pretocmd{\subsubsection}{\needspace{3\baselineskip}}{}{}

% Párrafos: penalizar líneas huérfanas (club) y viudas (widow) al cambiar de
% página: la primera línea de un párrafo no queda sola abajo, ni la última sola arriba.
\clubpenalty=10000
\widowpenalty=10000
\displaywidowpenalty=10000

% Matemáticas (variables como T, I, neq)
\usepackage{amsmath,amssymb}

% Imágenes: los diagramas PDF se ajustan al ancho de caja, sin anchos fijos
\usepackage{graphicx}
\makeatletter
\def\maxwidth{\ifdim\Gin@nat@width>\linewidth\linewidth\else\Gin@nat@width\fi}
\def\maxheight{\ifdim\Gin@nat@height>\textheight\textheight\else\Gin@nat@height\fi}
\makeatother
\setkeys{Gin}{width=\maxwidth,height=\maxheight,keepaspectratio}

% Tablas (los boletines pueden traer tablas Markdown)
\usepackage{longtable,booktabs,array,calc}

% Pies de figura
\usepackage{caption}
\captionsetup{font=small,labelfont=bf}

% Listas compactas de Markdown
\providecommand{\tightlist}{%
  \setlength{\itemsep}{0pt}\setlength{\parskip}{0pt}}

% Enlaces y destinos (al final, como recomienda hyperref)
\usepackage{hyperref}

\begin{document}

\hypertarget{termodinuxe1mica-financiera}{%
\section{TERMODINÁMICA FINANCIERA}\label{termodinuxe1mica-financiera}}

\textbf{Boletín TOC Insights --- Septiembre 2026}

\begin{center}\rule{0.5\linewidth}{0.5pt}\end{center}

\hypertarget{titular-y-resumen-ejecutivo}{%
\subsection{1. Titular y Resumen
Ejecutivo}\label{titular-y-resumen-ejecutivo}}

\textbf{El Cuello de Botella del Sistema Fiat: La Transición del
Throughput Global y la Saturación de Deuda Occidental}

Las presiones inflacionarias persistentes en EE. UU. (CPI en 3.4\%), el
fracaso de las intervenciones de liquidez del Tesoro (\$6,000M en
recompras), la escalada de los rendimientos europeos en máximos
multianuales y el inminente cambio de postura del Banco de Japón (18 de
septiembre) no son fenómenos aislados. Señalan la incapacidad del
mercado monetario mundial para absorber deuda sovereign sin respaldo en
Throughput (\(T\)) productivo real. La economía global ha iniciado un
desacoplamiento de la hegemonía del dólar estadounidense (similar a la
transición de la libra esterlina entre 1946 y 1948), redirigiendo las
líneas de producción hacia un sistema multipolar.

\begin{center}\rule{0.5\linewidth}{0.5pt}\end{center}

\hypertarget{diagnuxf3stico-del-sistema-seguxfan-toc}{%
\subsection{2. Diagnóstico del Sistema según
TOC}\label{diagnuxf3stico-del-sistema-seguxfan-toc}}

\hypertarget{la-restricciuxf3n-del-sistema}{%
\subsubsection{La Restricción del
Sistema}\label{la-restricciuxf3n-del-sistema}}

La capacidad finita del balance del mercado monetario mundial para
absorber emisiones masivas de deuda soberana de EE. UU. y Europa
(acumulación de \emph{Inventario Financiero \(I\)}) sin un crecimiento
equivalente en el \emph{Throughput (\(T\))} de la economía real
(generación neta de valor, bienes y energía).

\hypertarget{matriz-de-impacto-en-los-3-paruxe1metros-de-goldratt}{%
\subsubsection{Matriz de Impacto en los 3 Parámetros de
Goldratt}\label{matriz-de-impacto-en-los-3-paruxe1metros-de-goldratt}}

\begin{longtable}[]{@{}
  >{\raggedright\arraybackslash}p{(\columnwidth - 2\tabcolsep) * \real{0.5000}}
  >{\raggedright\arraybackslash}p{(\columnwidth - 2\tabcolsep) * \real{0.5000}}@{}}
\toprule\noalign{}
\begin{minipage}[b]{\linewidth}\raggedright
Parámetro TOC
\end{minipage} & \begin{minipage}[b]{\linewidth}\raggedright
Diagnóstico del Sistema Global
\end{minipage} \\
\midrule\noalign{}
\endhead
\bottomrule\noalign{}
\endlastfoot
\textbf{Throughput (\(T\))} & \textbf{Estancado / Desacoplado:} El flujo
nominal de crédito no genera bienes reales proporcionales. El \(T\)
físico migra fuera del circuito clearing en USD hacia transacciones
bilaterales y divisas alternativas (RMB, oro, swaps locales). \\
\textbf{Inventario (\(I\))} & \textbf{Saturado:} Acumulación desmedida
de títulos de deuda pública atascados en balances institucionales y
reservas de bancos centrales sin demanda privada suficiente al precio
nominal actual. \\
\textbf{Gasto de Operación (\(OE\))} & \textbf{Desbocado:} El costo del
servicio de deuda global se dispara a medida que los rendimientos a 10
años (4.85\%+ en EE. UU. y máximos en Europa) incrementan la carga
fiscal del sistema. \\
\end{longtable}

\begin{center}\rule{0.5\linewidth}{0.5pt}\end{center}

\hypertarget{herramientas-de-procesos-de-pensamiento-thinking-processes}{%
\subsubsection{Herramientas de Procesos de Pensamiento (Thinking
Processes)}\label{herramientas-de-procesos-de-pensamiento-thinking-processes}}

\hypertarget{uxe1rbol-de-realidad-actual-crt-saturaciuxf3n-de-deuda-y-contagio-global}{%
\paragraph{Árbol de Realidad Actual (CRT): Saturación de Deuda y
Contagio
Global}\label{uxe1rbol-de-realidad-actual-crt-saturaciuxf3n-de-deuda-y-contagio-global}}

\begin{longtable}[]{@{}
  >{\raggedright\arraybackslash}p{(\columnwidth - 4\tabcolsep) * \real{0.3333}}
  >{\raggedright\arraybackslash}p{(\columnwidth - 4\tabcolsep) * \real{0.3333}}
  >{\raggedright\arraybackslash}p{(\columnwidth - 4\tabcolsep) * \real{0.3333}}@{}}
\toprule\noalign{}
\begin{minipage}[b]{\linewidth}\raggedright
\end{minipage} & \begin{minipage}[b]{\linewidth}\raggedright
Causa Raíz / Restricción Fundamental
\end{minipage} & \begin{minipage}[b]{\linewidth}\raggedright
\end{minipage} \\
\midrule\noalign{}
\endhead
\bottomrule\noalign{}
\endlastfoot
\textbf{Efecto Indeseable 1 (UDE 1)} & \textbf{Efecto Indeseable 2 (UDE
2)} & \textbf{Efecto Indeseable 3 (UDE 3)} \\
Incapacidad de absorber deuda sovereign sin aumento de \(T\) real. &
Inyección de liquidez fiscal financiada con deuda del Tesoro. &
Rendimientos de deuda europea y estadounidense en máximos. \\
\textbf{Derivados de UDE 1} & \textbf{Derivados de UDE 2} &
\textbf{Derivados de UDE 3} \\
El \emph{Buyback} de \$6,000M destruye el precio de bonos viejos y eleva
yields. & El CPI repunta al 3.4\% interanual por distorsión de precios y
liquidez. & El BoJ se ve forzado a subir tasas (1.25\%), amenazando el
\emph{Yen Carry Trade}. \\
\textbf{Resultado Sistémico} & \textbf{Rediseño de Rutas de Flujo
(Bypass)} & \textbf{Fragmentación Multipolar} \\
Los tenedores extranjeros desinvierten de Bonos del Tesoro de EE. UU. &
Transacción de materias primas fuera del sistema SWIFT / USD. &
Reconfiguración de la Red de Cadena de Valor Global (RMB, Oro, Monedas
Locales). \\
\end{longtable}

\begin{center}\rule{0.5\linewidth}{0.5pt}\end{center}

\hypertarget{nube-de-evaporaciuxf3n-evaporating-cloud-el-conflicto-fiscal-vs.-monetario}{%
\paragraph{Nube de Evaporación (Evaporating Cloud): El Conflicto Fiscal
vs.~Monetario}\label{nube-de-evaporaciuxf3n-evaporating-cloud-el-conflicto-fiscal-vs.-monetario}}

\begin{longtable}[]{@{}
  >{\raggedright\arraybackslash}p{(\columnwidth - 4\tabcolsep) * \real{0.3333}}
  >{\raggedright\arraybackslash}p{(\columnwidth - 4\tabcolsep) * \real{0.3333}}
  >{\raggedright\arraybackslash}p{(\columnwidth - 4\tabcolsep) * \real{0.3333}}@{}}
\toprule\noalign{}
\begin{minipage}[b]{\linewidth}\raggedright
\end{minipage} & \begin{minipage}[b]{\linewidth}\raggedright
Objetivo Común (\(O\)): Estabilidad Económica Sistemática
\end{minipage} & \begin{minipage}[b]{\linewidth}\raggedright
\end{minipage} \\
\midrule\noalign{}
\endhead
\bottomrule\noalign{}
\endlastfoot
\textbf{Requisito A} & & \textbf{Requisito B} \\
Mantener la solvencia del mercado de deuda soberana y la liquidez del
Tesoro. & & Preservar el poder adquisitivo de la moneda y contener la
inflación. \\
\textbf{Prerrequisito A'} & & \textbf{Prerrequisito B'} \\
Recomprar deuda / monetizar el déficit (\emph{Quantitative Easing} o
\emph{Buybacks} masivos). & & Mantener tasas de interés elevadas y
restringir la hoja de balance del Banco Central. \\
& \textbf{Resolución del Conflicto (Inyección de Supuesto)} & \\
\textbf{Supuesto Falso:} El Dólar es la única divisa posible para
procesar el comercio global. & \textbf{Inyección TOC:} Diversificar el
comercio internacional fuera del clearing del USD (Transición
Multipolar), liberando la presión sobre el cuello de botella del balance
soberano de EE. UU. & \textbf{Falsación:} El flujo global de bienes
(\(T\)) ignora la restricción de liquidez en USD operando con vías de
pago alternativas. \\
\end{longtable}

\begin{center}\rule{0.5\linewidth}{0.5pt}\end{center}

\hypertarget{contabilidad-de-costos-tradicional-vs.-contabilidad-del-throughput-ta}{%
\subsection{3. Contabilidad de Costos Tradicional vs.~Contabilidad del
Throughput
(TA)}\label{contabilidad-de-costos-tradicional-vs.-contabilidad-del-throughput-ta}}

\begin{longtable}[]{@{}
  >{\raggedright\arraybackslash}p{(\columnwidth - 4\tabcolsep) * \real{0.3333}}
  >{\raggedright\arraybackslash}p{(\columnwidth - 4\tabcolsep) * \real{0.3333}}
  >{\raggedright\arraybackslash}p{(\columnwidth - 4\tabcolsep) * \real{0.3333}}@{}}
\toprule\noalign{}
\begin{minipage}[b]{\linewidth}\raggedright
Enfoque
\end{minipage} & \begin{minipage}[b]{\linewidth}\raggedright
Visión Ortodoxa / Contabilidad de Costos
\end{minipage} & \begin{minipage}[b]{\linewidth}\raggedright
Visión TOC / Contabilidad del Throughput
\end{minipage} \\
\midrule\noalign{}
\endhead
\bottomrule\noalign{}
\endlastfoot
\textbf{Interpretación del CPI y Yields} & Considera el repunte de la
inflación (3.4\%) como un desajuste temporal entre oferta y demanda
local que se arregla ajustando sutilmente la tasa de interés en 25 bps.
& Identifica la inflación como el síntoma de inyectar moneda sin
capacidad de Throughput productivo, elevando el Gasto de Operación
(\(OE\)) de toda la sociedad. \\
\textbf{Operaciones de Recompra (\emph{Buybacks})} & Ve las recompras de
deuda (\$6,000M) como intervenciones de ``eficiencia de
microestructura'' para inyectar liquidez y mejorar el funcionamiento del
mercado. & Lo define como la auto-compra de inventario defectuoso: el
Tesoro gasta recursos para adquirir sus propios pasivos sin demanda
real, aumentando el \(OE\) estatal sin resolver el cuello de botella. \\
\textbf{Solución Propuesta} & Ajustes locales de costo de endeudamiento,
recortes presupuestarios fragmentados y control de la velocidad del
balance del Banco Central. & Subordinación completa de la política
financiera a la capacidad real de producción. Rediseño de los canales de
distribución de valor. \\
\end{longtable}

\begin{center}\rule{0.5\linewidth}{0.5pt}\end{center}

\hypertarget{teoruxeda-monetaria-moderna-mmt-vs.-teoruxeda-de-restricciones-toc}{%
\subsection{4. Teoría Monetaria Moderna (MMT) vs.~Teoría de
Restricciones
(TOC)}\label{teoruxeda-monetaria-moderna-mmt-vs.-teoruxeda-de-restricciones-toc}}

\begin{figure}
\centering
\includegraphics{Diagramas/mmt-vs-toc.pdf}
\caption{Creación de dinero/liquidez: Visión MMT vs.~Visión TOC}
\end{figure}

\begin{itemize}
\tightlist
\item
  \textbf{Visión MMT:} Postula que el emisor soberano de moneda no
  enfrenta restricciones financieras nominales. Sostiene que el déficit
  público crea las reservas privadas y que el volumen de deuda no es un
  obstáculo mientras la inflación esté bajo control.
\item
  \textbf{Visión TOC:} Demuestra que \textbf{el dinero es únicamente un
  medio de transporte del flujo, no la restricción misma}. Si la emisión
  monetaria supera la capacidad de producción del cuello de botella real
  de la economía, el dinero extra no genera valor; simplemente se
  acumula como \emph{Inventario Financiero (\(I\))} tóxico y destruye la
  velocidad del \emph{Throughput (\(T\))} mediante inflación e
  incrementos de las tasas de interés.
\item
  \textbf{El Veredicto:} El escenario actual de 2026 invalida el
  supuesto central de la MMT: el mercado global ha dejado de absorber
  liquidez nominal no vinculada a bienes tangibles, castigando los
  títulos soberanos con yields del 4.85\%+ e imponiendo un límite físico
  infranqueable a la soberanía monetaria no respaldada por \(T\).
\end{itemize}

\begin{center}\rule{0.5\linewidth}{0.5pt}\end{center}

\hypertarget{conclusiuxf3n-y-prescripciuxf3n-estratuxe9gica-los-5-pasos-de-enfoque-de-goldratt}{%
\subsection{5. Conclusión y Prescripción Estratégica (Los 5 Pasos de
Enfoque de
Goldratt)}\label{conclusiuxf3n-y-prescripciuxf3n-estratuxe9gica-los-5-pasos-de-enfoque-de-goldratt}}

Para navegar la transición monetaria multipolar y la crisis de
saturación de deuda, la gestión macroeconómica y financiera debe aplicar
los \textbf{5 Pasos de Enfoque}:

\begin{enumerate}
\def\labelenumi{\arabic{enumi}.}
\tightlist
\item
  \textbf{Identificar la Restricción:} Reconocer que la restricción no
  es la falta de dinero nominal, sino la saturación del balance del
  mercado mundial para absorber deuda en USD sin un respaldo en
  Throughput real (\(T\)).
\item
  \textbf{Explotar la Restricción:} Maximizar el rendimiento real de la
  capacidad productiva actual sin emitir nueva deuda no respaldada.
  Orientar el gasto público exclusivamente a eliminar cuellos de botella
  de infraestructura, energía y tecnología.
\item
  \textbf{Subordinar todo lo demás a la Restricción:} Redireccionar las
  políticas del Banco Central y del Tesoro para alinearlas con el flujo
  real. Dejar de simular demanda mediante intervenciones cosméticas como
  \emph{buybacks} de \$6,000M que distorsionan el mercado de bonos.
\item
  \textbf{Elevar la Restricción:} Diversificar las plataformas de
  intercambio global. Adoptar acuerdos multilaterales de compensación en
  múltiples divisas (RMB, sistemas locales, respaldos en activos duros)
  para permitir que el Throughput del comercio global siga fluyendo sin
  pasar obligatoriamente por el cuello de botella de la deuda
  estadounidense.
\item
  \textbf{Si la Restricción se ha roto, Volver al Paso 1 (Prevenir la
  Inercia):} Una vez que el comercio mundial se desacople de la
  dependencia exclusiva del dólar, no permitir que la inercia cree un
  nuevo cuello de botella de sobre-emisión dentro de los nuevos bloques
  económicos multipolares.
\end{enumerate}

\end{document}
